\chapter{Зриви та електрони-утікачі}\label{ch:disruptions}

Зрив (disruption) --- це швидка втрата утримання: за мілісекунди плазма віддає стінці теплову енергію,
а струм плазми спадає. Електрони-утікачі (runaway electrons) --- це електрони, які поздовжнє поле
$\Epar$ прискорює швидше, ніж їх гальмують зіткнення. Вони з'являються і під час зриву, і на старті
розряду. Корпус описує утікачів докладніше, ніж сам зрив, тож зрив подано коротко, а основну увагу
приділено кінетиці утікачів~\cite{Hoppe2022,Marchuk2025}.

\section{Фази зриву}\label{sec:09-phases}

\begin{outofcorpus}[теплове і струмове гасіння, VDE, halo-струми]
\textbf{Теплове гасіння} (thermal quench). Після сильної МГД-нестійкості (часто після блокування моди
$2/1$) магнітні поверхні руйнуються, і теплова енергія виходить на стінку за $\sim0{,}1$--$1$~мс. $T_e$
падає від кеВ до одиниць--десятків еВ.
\textbf{Струмове гасіння} (current quench). Опір холодної плазми $\eta\propto T_e^{-3/2}$ зростає на
кілька порядків, і струм спадає за мілісекунди--десятки мілісекунд. Магнітний потік зникає, тому
індукується поле $\Epar\approx\Lpl|\dd\Ip/\dd t|/(2\pi R_0)$ ($\Lpl$~--- індуктивність плазми), на порядки більше за робоче
($\sim\SI{0,1}{В/м}$).
\textbf{VDE} (vertical displacement event). Витягнута плазма вертикально нестійка (розд.~\ref{ch:stability}).
Якщо керування втрачено, шнур зміщується вгору або вниз за час проникнення поля крізь камеру і
торкається стінки. VDE буває причиною зриву, а буває його наслідком.
\textbf{Halo-струми.} Коли шнур торкається стінки, частина полоїдального струму замикається через
відкриті силові лінії SOL і конструкції камери. Сили $\vect J\times\Bvec$ на камеру при цьому великі.
\end{outofcorpus}

У корпусі ці фази лише згадано. Код TSC обчислює halo-струм на відкритих лініях і його застосовували для сил зриву на пасивних стабілізаторах (PBX) і для моделювання VDE
(DIII-D). Струм, який з'являється в NSTX, коли до камери прикладено полоїдальну напругу (CHI), автори
порівнюють з halo-струмом зриву~\cite[с.~3--4]{Jardin2000}. Теплове гасіння в корпусі з'являється
лише один раз: як джерело «гарячого хвоста» (hot tail). Найшвидші електрони гарячого розподілу
гальмуються повільніше за основну масу і встигають розігнатися індукованим полем~\cite[с.~9]{Hoppe2022}.

\section{Генерація утікачів: механізми Драйсера і лавини}\label{sec:09-gen}

\begin{outofcorpus}[звідки беруться два порогові поля]
Сила тертя на швидкий електрон ($v\gg v_{Te}$) у головному порядку спадає як
$F(v)\approx e^4n_e\lnL/(4\pi\varepsilon_0^2m_ev^2)$ (множник порядку одиниці, що залежить від $\Zeff$,
тут опущено). При $v\to c$ вона виходить на мінімум $e\Ec$, тому поле $E<\Ec$ не може прискорити
жоден електрон. При $E>\Ec$ вільно прискорюються електрони з $v>v_c\approx c\sqrt{\Ec/E}$.
Для максвеллівського розподілу $v_c^2/v_{Te}^2=(m_ec^2/T_e)(\Ec/E)=\ED/E$, тому частка електронів у
«хвості» експоненційно мала за $\ED/E$.
\end{outofcorpus}

Корпус записує обидва пороги так~\cite[с.~9, (2.24)--(2.25)]{Hoppe2022}:
\begin{equation}\label{eq:09-Ec}
  \Ec=\frac{e^3n_e\lnL}{4\pi\varepsilon_0^2m_ec^2},\qquad
  \ED=\Ec\,\frac{m_ec^2}{T_e}=\frac{e^3n_e\lnL}{4\pi\varepsilon_0^2T_e}.
\end{equation}
Це поля Коннора--Гасті (Connor--Hastie) і Драйсера (Dreicer); у~\cite[с.~66--67, (3.2), (3.7)]{Marchuk2025}
їх записано в СГС (переведення~--- дод.~\ref{app:cgs}).
При $E\approx0{,}215\,\ED$ поле долає тертя для всіх електронів, і розподіл «зісковзує» (slide-away)
~\cite[с.~9]{Hoppe2022}. Обидва поля пропорційні $n_e$. Отже, за заданого $E$ розряд тим небезпечніший,
чим нижча густина: $E/\ED\propto T_e/n_e$.

Утікачі виникають двома шляхами. \emph{Первинна генерація} (Драйсер) --- це дифузійний
витік теплових електронів в область утікання. \emph{Лавина} --- це близькі зіткнення наявного утікача з
тепловим електроном, після яких новий електрон теж опиняється в області утікання. Баланс густини
утікачів~\cite[с.~10, (2.26)]{Hoppe2022}, \cite[с.~79, (3.14)]{Marchuk2025}:
\begin{equation}\label{eq:09-balance}
  \frac{\dd\nre}{\dd t}=\gD+\frac{\nre}{\tauav}-\frac{\nre}{\taure}.
\end{equation}
У~\cite{Marchuk2025} первинний член має вигляд $\gD=n_e\nu_e\lambda_e$~\cite[с.~69--70, (3.9)--(3.11)]{Marchuk2025},
де
\begin{equation}\label{eq:09-dreicer}
  \lambda_e=k_{\mathrm D}(\Zeff)\,\Big(\frac{\Epar}{\ED}\Big)^{-\frac{3(\Zeff+1)}{16}}
  \exp\!\Big[-\frac{\ED}{4\Epar}-\sqrt{\frac{(\Zeff+1)\ED}{\Epar}}\Big],\qquad
  \nu_e=\frac{e^4n_e\lnL}{4\pi\varepsilon_0^2m_e^2v_e^3}
\end{equation}
($\nu_e$ переведено з СГС; $k_{\mathrm D}(1)=0{,}32$, $k_{\mathrm D}(2)=0{,}43$). Показник $3(\Zeff+1)/16$
взято з робочої формули (3.26)~\cite[с.~84]{Marchuk2025}. У (3.10) на с.~69 множника 3 немає. Час
лавиноутворення (вигляд однаковий у СГС і СІ, бо входить лише $e\Epar$)~\cite[с.~71, (3.13)]{Marchuk2025}:
\begin{equation}\label{eq:09-tav}
  \tauav\approx\frac{\sqrt{12}\,m_ec\,\lnL\,(2+\Zeff)}{9\,e\Epar}.
\end{equation}
Модель~\cite{Marchuk2025} застосовна лише при $\Epar\ll\ED$~\cite[с.~66, (3.1)]{Marchuk2025}, і автор сам
називає це межею застосовності~\cite[с.~92]{Marchuk2025}. STREAM~\cite[с.~10]{Hoppe2022} обчислює
$\gD$ нейромережею Гесслова (Hesslow), а темп лавини --- напіваналітичною формулою з частковим
екрануванням ядер.

\section{Старт розряду: прогоряння і утікачі}\label{sec:09-start}

На старті густина низька. У ITER поле пробою обмежене значенням \SI{0,3}{В/м}~\cite[с.~2]{Hoppe2022}.
Поки триває прогоряння (burn-through), тобто нейтральний газ ще не іонізовано повністю, силові лінії ще відкриті.

\begin{corpusexample}[STREAM: ITER і JET \#77210~\cite{Hoppe2022}]
0D-код STREAM, побудований на DREAM, звірено з DYON на ідеалізованому сценарії ITER і на розряді JET
\#77210 з вуглецевою стінкою. Узгодження добре~\cite[с.~12--15]{Hoppe2022}. Параметри ITER: $p=\SI{0,8}{мПа}$,
$\Btor=\SI{2,65}{Тл}$, $R_0=\SI{5,65}{м}$, $a=\SI{1,6}{м}$, $\Vloop=\SI{12}{В}$; JET: \SI{2,7}{мПа},
\SI{2,4}{Тл}, $R_0=\SI{2,96}{м}$~\cite[с.~13, табл.~2]{Hoppe2022}.
\textbf{Базовий ITER} ($n_{\mathrm H}(0)=\SI{3,84e17}{м^{-3}}$~\cite[с.~12]{Hoppe2022}). Газ іонізується
за $\sim\SI{100}{мс}$. Попри $E/\ED$ до 85\,\% на старті, струм утікачів не виникає~\cite[с.~15--16]{Hoppe2022}:
поки магнітні поверхні відкриті (до $\Ip\approx\SI{100}{кА}$, $t\approx\SI{40}{мс}$~\cite[с.~17]{Hoppe2022}),
кожен новий утікач одразу йде на стінку, а коли поверхні замикаються, плазма вже настільки щільна й гаряча,
що $E/\ED$ практично нульове~\cite[с.~16]{Hoppe2022}. Зміна $\taure$ приблизно вдесятеро цього
не змінює~\cite[с.~21]{Hoppe2022}.
\textbf{Тиск у 10 разів нижчий} (\SI{0,08}{мПа}). Прогоряння проходить, проте омічна складова струму
зупиняється на $\approx\SI{1}{МА}$, а утікачі й далі набирають струм. Затравку весь час дає механізм
Драйсера; лавина стає помітною лише за кілька секунд~\cite[с.~16--17]{Hoppe2022}.
\textbf{Критерій небезпеки:} $E/\ED\gtrsim3\,\%$ протягом тривалого часу після прогоряння~\cite[с.~21]{Hoppe2022}.
\textbf{Напуск D} (\SI{2}{с}; старт через 0,5, 1 або \SI{3}{с} після пробою). Ранній старт (\SI{0,5}{с})
придушує майже всю генерацію Драйсера. Запізнілий (\SI{3}{с}) застає утікачів, що вже несуть $\sim75\,\%$
струму при $E/\ED\approx6\,\%$, і лавина далі розмножує цю затравку~\cite[с.~19--20]{Hoppe2022}. Проте
надто ранній чи надто інтенсивний напуск зриває прогоряння~\cite[с.~21]{Hoppe2022}. Автори наголошують,
що висновки якісні, а систему керування не враховано~\cite[с.~21]{Hoppe2022}.
\end{corpusexample}

\section{Струмове гасіння: 0D-модель сплеску поля}\label{sec:09-cq}

Під час зриву $E/\ED$ сягає помітних значень лише ненадовго, під час струмового гасіння. Потім поле
падає до ефективного порога, і далі утікачі розмножуються переважно лавиною~\cite[с.~17]{Hoppe2022}.
У моделі~\cite{Marchuk2025} сплеск задано кусково-лінійно. Поле зростає від $E_0$ до $E_0+E_{\mathrm m}$
за час $\tau_{\mathrm{br}}$ і за такий самий час повертається до $E_0$. $T_e(t)$ змінюється так само між
$T_1=\SI{325}{еВ}$ і $T_2=\SI{275}{еВ}$~\cite[с.~81--82, (3.16)--(3.17)]{Marchuk2025}. Утратами
знехтувано~\cite[с.~79]{Marchuk2025}. Тоді \eqref{eq:09-balance} --- лінійне рівняння першого
порядку, і його розв'язок має вигляд
\begin{equation}\label{eq:09-sol}
  \nre(t)=\mathcal A(t,t_0)\,\nre(t_0)+\int_{t_0}^{t}\mathcal A(t,t')\,\gD(t')\,\dd t',\qquad
  \mathcal A(t,t')=\exp\int_{t'}^{t}\Big(\frac1{\tauav}-\frac1{\taure}\Big)\dd t''
\end{equation}
($\mathcal A$~--- коефіцієнт підсилення; наш компактний запис формули (3.15)~\cite[с.~79--80]{Marchuk2025}). Після заміни змінних приріст
$\Delta\nre$ за сплеск зводиться до двох інтегралів, які обчислюють квадратурою (код RAEGh-np на Fortran)~\cite[с.~84--85]{Marchuk2025}. Квадратуру названо методом Гаусса на с.~85
і методом Сімпсона на с.~92.

\textbf{T-10}~\cite[с.~86--87]{Marchuk2025}: $n_e=\SI{0,9e19}{м^{-3}}$, $\Zeff=2$, $\lnL=15$, $E_0=\SI{0,8}{В/м}$,
$E_{\mathrm m}=0{,}5$--\SI{5}{В/м}, $2\tau_{\mathrm{br}}=\SI{0,25}{мс}$. Стрибок $\Delta\nre$ лежить у межах
$\sim10^{15}$--$10^{19}$\,м$^{-3}$ і зростає з $E_{\mathrm m}$ і $\tau_{\mathrm{br}}$. \emph{Твердження автора:}
у попередньому моделюванні T-10 (Savrukhin et al., 2015) вказано $\ED=\SI{35}{В/м}$ і
$E_{\mathrm m}=\SI{12,8}{В/м}$, тоді як насправді $\ED=\SI{13,36}{В/м}$ при \SI{275}{еВ} і \SI{11,32}{В/м} при \SI{325}{еВ}. Отже, умова $\Epar\ll\ED$ там
порушена~\cite[с.~74, 86]{Marchuk2025}. Наша оцінка за \eqref{eq:09-Ec}: \num{12,8} і \SI{10,8}{В/м}.
\textbf{EAST}~\cite[с.~76, 90]{Marchuk2025}: $R=\SI{1,86}{м}$, $\Ip=\SI{250}{кА}$, $\approx46\,\%$ струму
несуть надтеплові електрони. Модель спрощено, бо $T_e=\SI{550}{еВ}$ під час сплеску майже стала;
$n_e=\SI{2,2e19}{м^{-3}}$, $\Zeff=3$, $2\tau_{\mathrm{br}}=\SI{2,5}{мс}$. Узгодження з експериментом лише
якісне. Висновок автора: стабілізація МГД-активності зменшує амплітуду і тривалість сплесків $\Epar$
і тим послаблює генерацію утікачів~\cite[с.~91]{Marchuk2025}.

\section{Зв'язок із проєктом}\label{sec:09-project}

Прогноз зриву методами машинного навчання проєкт свідомо не взяв.

\begin{established}{E25}
Прогноз зриву \emph{насичений} і має \emph{закриті} дані: AUC $0{,}974$ на Alcator C-Mod в межах
однієї машини; DisruptionBench заархівовано без даних, дані JET мають юридичні обмеження~\cite{ProjRegisters}.
\end{established}
\noindent Відкритого симулятора зривів у проєкті теж немає. Рішення переглянуть, якщо
з'явиться відкритий датасет зривів із мітками моменту теплового гасіння.

\textbf{Відкриті питання проєкту} (Q4, Q5~\cite{ProjRegisters}). Це запитання до фахівців, не
гіпотези. Q4: які реалістичні стала часу зв'язування $\tau$ і запас
ентальпії $\Delta H$ надпровідників NbTi/Nb$_3$Sn у режимі ITER/JT-60SA? Q5: наскільки вакуумна камера
екранує швидке $\dd B/\dd t$ на котушках під час струмового гасіння? За екранування $\gtrsim10$ оцінки втрат
за моделлю Біна (Bean) і через сталу зв'язування втрачають сенс.

\subsection*{Контрольні запитання}
\begin{enumerate}
  \item Чому при $E<\Ec$ утікачів немає зовсім, хоча $\ED\gg\Ec$? Покажіть, що $\ED/\Ec=m_ec^2/T_e$.
  \item Чому в базовому сценарії ITER утікачів немає, хоча $E/\ED$ сягає 85\,\%? Чому зниження тиску
        в 10 разів це змінює~\cite{Hoppe2022}?
  \item Чому пізній напуск газу не зупиняє зростання струму утікачів, хоча придушує генерацію
        Драйсера?
  \item Чому модель \eqref{eq:09-balance}--\eqref{eq:09-tav} з~\cite{Marchuk2025} не можна застосувати до
        сплеску з $E_{\mathrm m}\approx\ED$?
\end{enumerate}

\subsection*{Задачі}
\begin{enumerate}
  \item \emph{Густина на старті ITER.} Візьміть $E=\SI{0,3}{В/м}$~\cite{Hoppe2022}, $\lnL=15$ і
        $n_e=\SI{3,84e17}{м^{-3}}$ (після повної іонізації). (а) Обчисліть $\Ec$ і $E/\Ec$. (б) За якої $T_e$ досягається поріг
        $E/\ED=3\,\%$? (в) Яка $n_e$ потрібна, щоб при $T_e=\SI{100}{еВ}$ виконувалося $E/\ED\le3\,\%$?
        Чи зникає при цьому лавина? Що зміниться при тиску в 10 разів нижчому?
  \item \emph{Сплеск на T-10.} За даними~\cite{Marchuk2025} ($n_e=\SI{0,9e19}{м^{-3}}$, $\lnL=15$, $\Zeff=2$)
        обчисліть $\ED$ при $T_e=275$ і \SI{325}{еВ}, а також $\Ec$. Знайдіть $\tauav$ за \eqref{eq:09-tav}
        при $\Epar=0{,}8$, 5 і \SI{12,8}{В/м}. Оцініть, у скільки разів лавина підсилює $\nre$ за сплеск
        $2\tau_{\mathrm{br}}=\SI{0,25}{мс}$ з $E_{\mathrm m}=\SI{5}{В/м}$. Що визначає стрибок $\Delta\nre$?
\end{enumerate}
